\subsection{Why diffusion capacity matters}
\label{subsec:diffusion_capacity}

Diffusion capacity is the economic margin that determines whether frontier AI capability becomes productive capacity. Frontier capability describes the ability to build models, assemble compute, produce algorithms, and advance technical performance. Diffusion capacity describes the ability to reorganize firms, sectors, cities, and production systems around those capabilities. This distinction matters because the economic effects of a general-purpose technology emerge through complements that are expensive, local, organizational, and slow to build. AI can expand the technological frontier while measured productivity remains limited if firms lack the data systems, managerial routines, equipment interfaces, software integration, and workforce capabilities needed to use it.

The adoption margin is especially important for artificial intelligence because AI often enters the economy through tasks and workflows rather than through a single observable capital good. In manufacturing, AI may appear as machine vision, predictive maintenance, scheduling optimization, robotics coordination, automated inspection, process control, warehouse management, or digital-twin simulation. These applications are embedded inside production systems, which makes them harder to observe than model releases, chip purchases, or venture investment. A country can therefore appear strong in frontier AI indicators while remaining weak in industrial absorption, and another country can accumulate strategic advantage through deployment channels that are largely invisible to frontier rankings.

The diffusion-capacity perspective changes how AI competition should be measured. Conventional capability indicators ask which actors can produce powerful systems. A diffusion-capacity measure asks which institutions can translate systems into output, quality, reliability, throughput, and industrial coordination. The second question requires attention to municipal governments, standards, industrial parks, procurement systems, demonstration projects, technical service providers, sectoral production routines, and administrative recognition mechanisms. These institutions shape the cost of adoption, the credibility of implementation, the circulation of technical knowledge, and the ability of firms to coordinate complementary investments.

China provides a useful empirical setting because its industrial AI strategy makes several layers of diffusion capacity visible. National AI pilot zones identify places expected to organize AI application and demonstration. Excellence-level smart-factory recognition identifies projects that have entered a formal category of intelligent manufacturing. City-level industrial ecosystems provide the firms, suppliers, engineers, and local policy capacity needed to implement production technologies. Sectoral composition determines where AI-enabled production systems are technically plausible. The paper uses this institutional visibility to measure a diffusion architecture that would be difficult to observe through firm-level AI-use surveys alone.

Diffusion capacity also explains why the paper focuses on hubs rather than treating policy designation as a flat treatment. General-purpose technologies often diffuse through locations where complementary assets are already concentrated. Large hubs contain engineering labor, supplier networks, universities, leading firms, industrial parks, capital, logistics systems, and administrative experience. These assets reduce the cost of implementation and increase the likelihood that a recognized project can be planned, executed, and reported. A positive association between pilot-zone designation and smart-factory recognition therefore carries the most economic meaning when interpreted through the geography of complementary capacity.

The concept of diffusion capacity turns the paper's empirical design into a measurement contribution. Smart-factory recognition is used because it provides a visible administrative trace of production-side adoption. City-resolution evidence is audited because the relevant adoption architecture is spatial. Hub exclusions are central because they reveal whether the association is carried by high-capacity locations. Sectoral exposure is measured because AI adoption depends on compatibility between technology and production processes. Export relevance is reported because diffusion in strategically important sectors has implications for industrial competition. The resulting design measures how AI becomes embedded in production, rather than only where AI capability is produced.

The diffusion-capacity framework also disciplines the paper's causal interpretation. The evidence can show that pilot-zone designation, smart-factory recognition, hub concentration, sectoral compatibility, and export-relevant manufacturing align in a coherent adoption architecture. The evidence cannot by itself show that designation caused the projects, that recognition raised productivity, or that observed sectors experienced export upgrading. The contribution lies in making an otherwise hidden margin of AI competition observable with a claim hierarchy that distinguishes architecture, association, mechanism evidence, robustness, and causal identification.